\section{Conclusion and Limitations}
\label{sec:conclusion}

We specify a pose-driven MRAC counter for people whose pace changes within a
track. The central design separates identity state while sharing parameters,
routes overlapping windows from local period evidence, fuses continuous
responses by normalized overlap-add, and decodes each track once. The setting
does not remove all annotation and does not jointly optimize perception and
counting: the counting objective is intended to avoid person-wise count and
period labels, while upstream perception may use external supervision.

This draft has three material limitations. The current local checkout contains
only a single-person PAMS reproduction; the multi-person data path, router,
continuity loss, and boundary decoder await synchronization with collaborator
code. No eligible MultiRep artifact has been received, so the manuscript makes
no performance claim. Finally, a modular pose/track pipeline inherits missed
detections and identity switches and may scale with the number of active
people. The final study must quantify each limitation through oracle/predicted
tracks, corruption analysis, and population-scaling measurements. Until those
gates close, the paper remains provisional and data-pending. MultiRep is
synthetic, and single-person transfer cannot establish multi-person identity
preservation; absent a real multi-person evaluation, any eventual conclusion
must remain restricted to MultiRep and the declared controlled tracker
corruptions.
